\documentclass[10pt,a4paper]{amsart}
\usepackage[margin=19mm]{geometry}
\usepackage{amsmath,amssymb,mathtools}
\usepackage[scr=rsfs,cal=euler]{mathalfa}
\usepackage{xfrac}
\usepackage[unicode,hidelinks,bookmarks=false]{hyperref}
\newcommand{\ie}{i.e.\ }
\newcommand{\cf}{cf.\ }
\newcommand{\resp}{respectively\ }
\newcommand{\Subsection}{\subsection}
\providecommand{\nobreakdash}{\nobreak\hbox{-}}
\setlength{\emergencystretch}{2em}
\multlinegap0pt
\usepackage{amssymb}
\usepackage{enumerate}
\usepackage{bm}
\newtheorem{theorem}{Theorem}
\def \De {\Delta}
\def\Om{\Omega}
\def \a {\alpha }
\def\pa{\partial}
\def\w {\omega}
\title{$L^{2}$-Hodge theory on Complete Almost K\"{a}hler Manifolds and the Hopf Conjecture}
\author{Teng Huang, Qiang Tan, Pan Zhang}
\date{Source: arXiv:2302.14032v4 (CC BY 4.0)}
\begin{document}
\maketitle

\noindent\emph{Source and licence.} $L^{2}$-Hodge theory on Complete Almost K\"{a}hler Manifolds and the Hopf Conjecture. Version arXiv:2302.14032v4 (CC BY 4.0), distributed by arXiv under the Creative Commons Attribution 4.0 International licence, http://creativecommons.org/licenses/by/4.0/, as recorded in the arXiv licence metadata for this exact version and shown on the abstract page https://arxiv.org/abs/2302.14032v4. Full licence notice: \texttt{licenses/arxiv-2302.14032-CC-BY-4.0.txt}.\par\smallskip

\noindent\textit{Editorial scope.} Counted unit: the article's theorem Thm3 (source0.tex lines 1022--1030). The article's complete original proof of it is delivered with it (source0.tex lines 1032--1087, one \texttt{proof} environment of 56 lines). No mathematics is written, repaired or re-derived: the statement and the proof are the article's own lines, in the article's own notation, and no retained line is substituted or re-flowed.  Retained uncounted as the source-local context the proof needs: lines 312--322; lines 1005--1017.  Nothing was omitted from any retained span, so the package carries no omission record.  No retained line contains a citation, so the wrapper carries no bibliography and no bibliography entry.  No retained line contains a figure environment, an image-inclusion command or drawing code, so no image is delivered and the package declares no asset dependency.  Editorial formatting only, in addition: an AMS \texttt{amsart} preamble that restates the article's own declarations and macros used by the retained lines, namely the environments \texttt{theorem} on separate counters, together with the standard packages those lines need.  The retained environments print in this document as Theorem 1 in order of appearance, whereas the article prints the same items with the numbering of its own full document.  The complete native source of the article as distributed in the arXiv e-print for this exact version is bundled unchanged as \texttt{sources/137931/source0.tex}.
	\begin{equation}\label{E27}
	\begin{split}
	&\mu^{2}=0,\\
	&\mu\pa+\pa\mu=0,\\
	&\pa^{2}+\mu\bar{\pa}+\bar{\pa}\mu=0,\\
	&\pa\bar{\pa}+\bar{\pa}\pa+\mu\bar{\mu}+\bar{\mu}\mu=0,\\
	&\bar{\pa}^{2}+\bar{\mu}\pa+\pa\bar{\mu}=0,\\
	&\bar{\mu}\bar{\pa}+\bar{\pa}\bar{\mu}=0,\\
	&\bar{\mu}^{2}=0.\\
	\end{split}
	\end{equation} 

	\begin{theorem}\label{Thm6}
	Let $(X,J,\w)$ be a complete $2n$-dimensional almost K\"{a}hler manifold. Suppose that there exists a bounded $1$-form $\theta$ such that
	\[\w=d\theta.\]
	Then for any $\a\in\Om^{p,q}_{0}(X)$ with $k:=p+q\neq n$, the following estimates hold:
	\begin{equation*}
	\begin{split}
	&\|\a\|_{L^{2}(X)}\leq c_2(n,k)\|\theta\|_{L^{\infty}(X)}\Big((\De_{\pa}+\De_{\bar{\pa}})\a,\a\Big)^{\frac{1}{2}},\\
	&\|\a\|_{L^{2}(X)}^{2}\Big(1-c_2(n,k)^{2}\|\theta\|^{2}_{L^{\infty}(X)}\sup|N_{J}|^{2}\Big)
	\leq 2c_2(n,k)^{2}\|\theta\|^{2}_{L^{\infty}(X)}(\De_{\bullet}\a,\a),
	\end{split}
	\end{equation*}
	where $\bullet=\pa,\bar{\pa}$. 
\end{theorem}	

\begin{theorem}\label{Thm3}
	Let $(X,J,\w)$ be a complete $2n$-dimensional almost K\"{a}hler manifold. Suppose that there exists a bounded $1$-form $\theta$ such that
	\[ \w=d\theta.\]
	Then there exist constants $c_{3}(n),c_{4}(n)$ depending only on $n$ such that the following holds: 
	
	for any $\a\in\Om_{0}^{p,+}(X)$ when $n-p$ is odd (resp. $\a\in\Om_{0}^{p,-}(X)$ when $n-p$ is even), 
	\[\|\theta\|^{2}_{L^{\infty}(X)}\|(\bar{\partial}+\bar{\partial}^{\ast})\a\|_{L^{2}(X)}
	\geq \Big(c_{3}(n)-c_{4}(n)\|\theta\|^{2}_{L^{\infty}(X)}\sup|N_{J}|^{2}\Big)\|\a\|^{2}_{L^{2}(X)}. \]
\end{theorem}

\begin{proof}
	We give the proof for the case when $n-p$ is odd; the even case is entirely analogous. 
	
	First, observe that
	\[({\bar{\partial}+\bar{\partial}^{\ast}})^{2}=\De_{\bar{\pa}}+(\bar{\pa})^{2}+(\bar{\pa}^{\ast})^{2}. \]
	For any $\a\in\Om_{0}^{p,+}(X)$, we decompose  $\a:=\sum\a_{2i}$ with $\a_{2i}\in\Om^{p,2i}$. Then
	\begin{equation*}
	\begin{split}
	\|({\bar{\partial}+\bar{\partial}^{\ast}})\a\|_{L^{2}(X)}^{2}
	&=(\De_{\bar{\pa}}\a,\a)+\Big((\bar{\pa})^{2}\a,\a\Big)+\Big((\bar{\pa}^{\ast})^2\a,\a\Big)\\
	&=(\De_{\bar{\pa}}\a,\a)+2\mathrm{Re}\Big((\bar{\pa})^{2}\a,\a\Big)\\
	&=\sum(\De_{\bar{\pa}}\a_{2i},\a_{2i})+2\sum\mathrm{Re}\Big((\bar{\pa})^{2}\a_{2i},\a_{2i+2}\Big)\\
	&=\sum(\De_{\bar{\pa}}\a_{2i},\a_{2i})-2\sum\mathrm{Re}\Big([\pa,\bar{\mu}]\a_{2i},\a_{2i+2}\Big),
	\end{split}
	\end{equation*}
	where we used the identity (cf.(\ref{E27}))
	\[ \bar{\pa}^{2}=-[\pa,\bar{\mu}].\]
	
	For each term in the sum, we estimate:
	\begin{equation*}
	\begin{split}
	|([\pa,\bar{\mu}]\a_{2i},\a_{2i+2})|&=|(\bar{\mu}\a_{2i},\pa^{\ast}\a_{2i+2})+(\pa\a_{2i},\bar{\mu}^{\ast}\a_{2i+2})|\\
	&\leq \|\bar{\mu}\a_{2i}\|_{L^{2}}\|\pa^{\ast}\a_{2i+2}\|_{L^{2}}+ \|\pa\a_{2i}\|_{L^{2}}\|\bar{\mu}^{\ast}\a_{2i+2}\|_{L^{2}}.\\
	\end{split}
	\end{equation*}
	Summing over $i$ and applying the elementary inequality $2ab\leq c_{0}a^{2}+\frac{1}{c_{0}}b^{2}$
	for all $c_{0}>0$, we obtain
	\begin{equation*}
	\begin{split}
	|2\sum\mathrm{Re}([\pa,\bar{\mu}]\a_{2i},\a_{2i+2})|
	&\leq2\sum(\|\bar{\mu}\a_{2i}\|_{L^{2}}\|\pa^{\ast}\a_{2i+2}\|_{L^{2}}+ \|\pa\a_{2i}\|_{L^{2}}\|\bar{\mu}^{\ast}\a_{2i+2}\|_{L^{2}})\\
	&\leq \sum c_{0}(\|\pa\a_{2i}\|^2_{L^{2}}+\|\pa^{\ast}\a_{2i}\|^2_{L^{2}})
	+\sum\frac{1}{c_{0}}(\|\bar{\mu}\a_{2i}\|^2_{L^{2}}+\|\bar{\mu}^{\ast}\a_{2i}\|^2_{L^{2}})\\
	&=\sum c_{0}(\De_{\pa}\a_{2i},\a_{2i})+\sum\frac{1}{c_{0}}(\De_{\bar{\mu}}\a_{2i},\a_{2i}).\\
	\end{split}
	\end{equation*}	
	Combining this with the earlier expression for $\|({\bar{\partial}+\bar{\partial}^{\ast}})\a\|^{2}$ and using the identity $$\De_{\bar{\pa}}=\De_{\pa}+\De_{\bar{\mu}}-\De_{\mu},$$
	we find 
	\begin{equation*}
	\begin{split}
	\|({\bar{\partial}+\bar{\partial}^{\ast}})\a\|_{L^{2}(X)}^{2}&\geq\sum\big{(}(\De_{\bar{\pa}}\a_{2i},\a_{2i})- c_{0}(\De_{\pa}\a_{2i},\a_{2i})-\frac{1}{c_{0}}(\De_{\bar{\mu}}\a_{2i},\a_{2i})\big{)}\\
	&=\sum\big{(}(1-c_{0})(\De_{\bar{\pa}}\a_{2i},\a_{2i})- c_{0}(\De_{\mu}\a_{2i},\a_{2i})-(\frac{1}{c_{0}}-c_{0})(\De_{\bar{\mu}}\a_{2i},\a_{2i})\big{)}.\\
	\end{split}
	\end{equation*}
	
	Choosing $c_{0} = \frac{1}{2}$ and applying Theorem \ref{Thm6}, we obtain
	\begin{equation*}
	\begin{split}
	\|({\bar{\partial}+\bar{\partial}^{\ast}})\a\|_{L^{2}(X)}^{2}
	&\geq\sum\Big(\frac{1}{2}(\De_{\bar{\pa}}\a_{2i},\a_{2i})-\frac{1}{2}(\De_{\mu}\a_{2i},\a_{2i})-\frac{3}{2}(\De_{\bar{\mu}}\a_{2i},\a_{2i})\Big)\\
	&\geq \sum \|\theta\|^{-2}_{L^{\infty}(X)}\Big(c_{3}(n)-c_{4}(n)\|\theta\|^{2}_{L^{\infty}(X)}\sup|N_J|^{2}\Big)\|\a_{2i}\|^{2}_{L^{2}(X)}\\
	&= \|\theta\|^{-2}_{L^{\infty}(X)}\Big(c_{3}(n)-c_{4}(n)\|\theta\|^{2}_{L^{\infty}(X)}\sup|N_J|^{2}\Big)\|\a\|_{L^{2}(X)}^{2}.
	\end{split}
	\end{equation*}
	This completes the proof.
\end{proof}

\end{document}
